% Insertion-ready research notes; uses standard amsmath/amsthm notation.
% Derivation and numerical audit: 2026-10-10.
% Do not advertise a new general parity principle: see the scope paragraph.

\section{A two-shift Lerch reflection for inverse-color doubles}
\label{gap:sec:main}\label{sec:gap}

The manuscript already proves an all-weight reduction of
$\operatorname{Li}_{a,b}(z,z^{-1})$ to one-direction double
polylogarithms. Its explicit mixed parity formula treats $b=1$; the
general parity component also follows by combining that reduction with
the all-index one-direction parity theorem. Thus an extension from
$b=1$ alone must not be advertised as a new reducibility principle.
The established multiple-polylogarithm parity theory includes Panzer's
general theorem and explicit depth-two formulas \cite{panzer}. More recently, Xu's
2026 preprint \emph{Symmetry Results for Cyclotomic Multiple Hurwitz
Zeta Values via Contour Integrals}, arXiv:2602.10391 \cite{xu}, proves broad
shifted cyclotomic symmetry identities. Its index order is increasing,
whereas the manuscript uses decreasing indices.

The purpose of the following calculation is more specific: it gives a
compact two-shift identity that is an ordinary holomorphic identity in
the full disk of a continuous color variable. Its arbitrary-index
version follows from a single first-order seed by parameter
differentiation. It also gives complete reductions, rather than only
paired reductions, on the complementary-shift diagonal. These are
useful explicit additions to the manuscript. No claim of worldwide
priority is made for the underlying symmetry.

\subsection{Definitions and analytic domain}

For integers $a,b\geq1$, define
\begin{equation}\label{gap:eq:def}
 \begin{aligned}
 \mathcal M_{a,b}(z;u,v)
 &=\sum_{n>m\geq0}
   \frac{z^{n-m}}{(n+u)^a(m+v)^b}
 =\sum_{h\geq1}z^h c_{a,b}(h;u,v),\\
 c_{a,b}(h;u,v)&=\sum_{m\geq0}
       \frac1{(m+h+u)^a(m+v)^b}.
 \end{aligned}
\end{equation}
We initially work on
\begin{equation}\label{gap:eq:domain}
 |z|<1,\qquad 0<\operatorname{Re}u<1,
 \qquad 0<\operatorname{Re}v<1.
\end{equation}
All powers are integer powers, so there is no denominator branch
choice. The Lerch function is
\[
 \Phi(z,s,q)=\sum_{m\geq0}\frac{z^m}{(m+q)^s}
 \quad (|z|<1,\ \operatorname{Re}q>0).
\]
Put
\begin{equation}\label{gap:eq:C}
 C_1(q)=\pi\cot\pi q,
 \qquad
 C_r(q)=\zeta(r,q)+(-1)^r\zeta(r,1-q)\quad(r\geq2).
\end{equation}
Equivalently,
\begin{equation}\label{gap:eq:Cderivative}
 C_r(q)=\frac{(-1)^{r-1}}{(r-1)!}
          \frac{d^{r-1}}{dq^{r-1}}(\pi\cot\pi q).
\end{equation}
Thus every $C_r$ is elementary in trigonometric functions, and can
equally be expressed through polygamma functions.

The series in \eqref{gap:eq:def} converges absolutely and normally on
compact subsets of \eqref{gap:eq:domain}. Indeed, uniformly on compact
shift sets,
\[
 |c_{a,b}(h;u,v)|\leq
 C\sum_{m\geq0}\frac1{(m+h+1)^a(m+1)^b}
 =O\left(\frac{1+\log(h+1)}{(h+1)^a}\right).
\]
The same estimates hold after any fixed number of shift derivatives.
In particular all parameter differentiations below are legitimate
holomorphic differentiations, not formal manipulations of a divergent
series.

\subsection{The first-order seed}

\begin{theorem}[Two-shift seed]\label{gap:thm:seed}
On \eqref{gap:eq:domain},
\begin{align}
 &\mathcal M_{1,1}(z;u,v)
   +\mathcal M_{1,1}(z;1-v,1-u)\notag\\
 &\quad=z\Phi(z,1,u)\Phi(z,1,1-v)
   +\bigl(C_1(v)-C_1(u)\bigr)
       z\Phi(z,1,1+u-v).
 \label{gap:eq:seed}
\end{align}
All apparent confluences at $u=v$ are removable; in fact no divided
difference is needed in the displayed formula.
\end{theorem}

\begin{proof}
Fix a positive integer $h$. The absolutely convergent bilateral sum
\[
 B_{1,1}(h;u,v)=\sum_{m\in\mathbb Z}
       \frac1{(m+h+u)(m+v)}
\]
has three disjoint ranges. The range $m\geq0$ is $c_{1,1}(h;u,v)$.
In the range $m\leq-h-1$, set $m=-h-1-\ell$; its contribution is
$c_{1,1}(h;1-v,1-u)$. In the middle range $-h\leq m\leq-1$, write
$m=-\ell-1$ and $j=h-1-\ell$. Its contribution is
\[
 -\sum_{j+\ell=h-1}\frac1{(j+u)(\ell+1-v)}.
\]
On the other hand, with $d=h+u-v$,
\[
 \frac1{(m+h+u)(m+v)}
 =\frac1d\left(\frac1{m+v}-\frac1{m+h+u}\right).
\]
Symmetric summation gives
$B_{1,1}(h;u,v)=(C_1(v)-C_1(u))/(h+u-v)$.
The shift by the integer $h$ does not change a symmetric cotangent
sum: the finitely shifted endpoint terms tend to zero. Multiply by
$z^h$ and sum over $h\geq1$. The middle range is the coefficient
convolution of the two Lerch series, and the bilateral term sums to
the second term in \eqref{gap:eq:seed}. All resulting $h$-series
converge absolutely for $|z|<1$.
\end{proof}

\subsection{Every pair of indices, in two equivalent forms}

\begin{theorem}[Arbitrary-index reflection]\label{gap:thm:all}
Let $a,b\geq1$ and $w=a+b$. Set $q=1+u-v$. Then
\begin{align}
 &\mathcal M_{a,b}(z;u,v)
      +(-1)^w\mathcal M_{b,a}(z;1-v,1-u)\notag\\
 &=z\Bigg[
 (-1)^b\sum_{r=1}^{a}\binom{w-r-1}{b-1}
       C_r(u)\Phi(z,w-r,q)\notag\\
 &\hspace{17mm}
 +\sum_{r=1}^{b}(-1)^{b-r}\binom{w-r-1}{a-1}
       C_r(v)\Phi(z,w-r,q)\notag\\
 &\hspace{17mm}
 -(-1)^b\Phi(z,a,u)\Phi(z,b,1-v)\Bigg].
 \label{gap:eq:all}
\end{align}
Equivalently, the right side is
\begin{equation}\label{gap:eq:all-derivative}
 \frac{(-1)^w}{(a-1)!(b-1)!}
 \partial_u^{a-1}\partial_v^{b-1}
 \left[
 z\Phi(z,1,u)\Phi(z,1,1-v)
 +(C_1(v)-C_1(u))z\Phi(z,1,1+u-v)
 \right].
\end{equation}
\end{theorem}

\begin{proof}
For the derivative form, apply the stated mixed derivative to
\eqref{gap:eq:seed}. On the first summand on its left, the derivative
contributes the sign $(-1)^{a+b-2}=(-1)^w$ and the factorial factor.
On the reflected summand both chain-rule signs cancel the signs
from differentiating a reciprocal, and the two indices are exchanged.
This gives \eqref{gap:eq:all-derivative} with the precise sign shown.

For an independent coefficient proof of \eqref{gap:eq:all}, put
$d=h+u-v$. The rational partial-fraction identity is
\begin{align}
 \frac1{(m+h+u)^a(m+v)^b}
 ={}&(-1)^b\sum_{r=1}^{a}
  \binom{w-r-1}{b-1}\frac1{d^{w-r}(m+h+u)^r}\notag\\
 &+\sum_{r=1}^{b}(-1)^{b-r}
  \binom{w-r-1}{a-1}\frac1{d^{w-r}(m+v)^r}.
 \label{gap:eq:pf}
\end{align}
It follows either by expanding at its two poles or by subtracting
the two sides, observing that the difference has no finite poles and
vanishes at infinity. Since $\operatorname{Re}d>0$, no collision occurs
in the domain under consideration.

Sum over $m\in\mathbb Z$. For $r\geq2$, the sum of each reciprocal
power is $C_r(u)$ or $C_r(v)$. The $r=1$ terms are evaluated by a
common symmetric cutoff. Their principal parts at infinity cancel,
because
\[
 \binom{w-2}{a-1}=\binom{w-2}{b-1}.
\]
The original bilateral sum is absolutely convergent, so this use of
a common cutoff changes no value. Split it into the same three ranges
as in the first-order proof. The negative range is now
$(-1)^w c_{b,a}(h;1-v,1-u)$, and the middle range is
\[
 (-1)^b\sum_{j+\ell=h-1}
      (j+u)^{-a}(\ell+1-v)^{-b}.
\]
Multiplication by $z^h$ and summation gives
\eqref{gap:eq:all}.
\end{proof}

\subsection{Shift derivatives and normalized primitives}

The whole family is closed under arbitrary shift derivatives:
\begin{equation}\label{gap:eq:shift-jets}
 \partial_u^r\partial_v^s\mathcal M_{a,b}(z;u,v)
 =(-1)^{r+s}(a)_r(b)_s\mathcal M_{a+r,b+s}(z;u,v).
\end{equation}
Accordingly, for $a\geq2$ a primitive in $u$ is
$-\mathcal M_{a-1,b}/(a-1)$, and for $b\geq2$ a primitive in $v$ is
$-\mathcal M_{a,b-1}/(b-1)$. At the index-one boundary one should
normalize the primitive before summing. For a path between $u_0$ and
$u_1$ inside the right half-plane,
\begin{equation}\label{gap:eq:log-primitive}
 \int_{u_0}^{u_1}\mathcal M_{1,b}(z;u,v)\,du
 =\sum_{n>m\geq0}\frac{z^{n-m}}{(m+v)^b}
       \log\frac{n+u_1}{n+u_0}.
\end{equation}
Here the logarithm is the difference of the common analytic logarithms
along the path. For $|z|<1$ the normalized series is absolutely
convergent, including $b=1$, because the logarithmic difference is
$O(n^{-1})$ on compact endpoint sets. The unnormalized series with
$\log(n+u)$ alone need not converge when $b=1$.

There is also a precise spectral-derivative interpretation. Define the
termwise difference
\[
 \Delta\mathcal M_{s,b}
 =\sum_{n>m\geq0}\frac{z^{n-m}}{(m+v)^b}
       \bigl((n+u_1)^{-s}-(n+u_0)^{-s}\bigr).
\]
It is holomorphic near $s=0$, even when the two separate sums there
are divergent, and the right side of \eqref{gap:eq:log-primitive} is
$-\partial_s\Delta\mathcal M_{s,b}|_{s=0}$. This statement uses
subtraction before continuation. It does not assert that the
integer-index finite sum \eqref{gap:eq:all} can itself be
differentiated in a noninteger exponent.

\paragraph{Boundary values.}
The identity extends to each $|z|=1$, $z\ne1$, as an ordinary
convergent identity. To make this statement independent of an
unspecified Abel prescription, note that on compact shift sets
$c_{a,b}(h;u,v)\to0$ and
\[
 c_{a,b}(h+1;u,v)-c_{a,b}(h;u,v)
 =O\left(\frac{1+\log(h+1)}{(h+1)^{a+1}}\right).
\]
The sum of these differences converges absolutely. Summation by
parts therefore proves convergence uniformly on closed arcs that
avoid $1$. The same argument applies to every Lerch term on the
right, whose parameter has positive real part, and to fixed shift
derivatives. Taking a radial limit of the disk identity proves the
boundary identity in the ordinary gap-ordered series convention.

It also agrees with the original outer-index order. Let $O_N$ be the
sum in \eqref{gap:eq:def} restricted to $n\leq N$, and let
$D_N=\sum_{h=1}^Nz^hc_{a,b}(h;u,v)$. Partial summation of the finite
$h$-sum gives, uniformly for shifts in a compact subset of the strip,
\[
 |D_N-O_N|\leq\frac{C}{|1-z|}
 \left(N^{-a}(1+\log(N+1))+\sum_{m>N}m^{-a-b}\right)
 \longrightarrow0.
\]
For $m\leq N$, the omitted $h$-range starts past $N-m$, so its
denominator $(m+h+u)^a$ is of size at least $N^a$; for $m>N$, the
full finite $h$-sum is bounded by a constant times $m^{-a}$.
These are exactly the two terms in the displayed bound. Thus no
conditional rearrangement is hidden in the boundary statement.
If $a,b\geq2$,
all terms are also absolutely convergent at $z=1$, so that endpoint
may then be included.

\subsection{Complete reductions at complementary shifts}

The reflection becomes a full evaluation when the two exponents
agree and $v=1-u$.

\begin{corollary}[Complementary diagonal]\label{gap:cor:diagonal}
For $a\geq1$, $|z|<1$ and $0<\operatorname{Re}u<1$,
\begin{align}
 \mathcal M_{a,a}(z;u,1-u)
 =(-1)^a z\Bigg[
 &\sum_{r=1}^{a}\binom{2a-r-1}{a-1}
        C_r(u)\Phi(z,2a-r,2u)\notag\\
 &-\frac12\Phi(z,a,u)^2\Bigg].
 \label{gap:eq:diagonal}
\end{align}
In particular,
\begin{equation}\label{gap:eq:diagonal11}
 \mathcal M_{1,1}(z;u,1-u)
 =\frac z2\Phi(z,1,u)^2
   -\pi\cot(\pi u)\,z\Phi(z,1,2u),
\end{equation}
and
\begin{align}
 \mathcal M_{2,2}(z;u,1-u)
 =z\bigl[&2\pi\cot(\pi u)\Phi(z,3,2u)
          +\pi^2\csc^2(\pi u)\Phi(z,2,2u)\notag\\
          &-\tfrac12\Phi(z,2,u)^2\bigr].
 \label{gap:eq:diagonal22}
\end{align}
\end{corollary}
\begin{proof}
The reflected double on the left of \eqref{gap:eq:all} equals the
original one. On the right, use $C_r(1-u)=(-1)^rC_r(u)$ and divide
by two.
\end{proof}

The $a=1$ identity also exposes the generalized-harmonic content.
With
\[
 H_h(u)=\psi(h+u)-\psi(u)=\sum_{j=0}^{h-1}\frac1{j+u},
\]
partial fractions give
\[
 c_{1,1}(h;u,1-u)
 =\frac{H_h(u)-\pi\cot\pi u}{h+2u-1}.
\]
Consequently \eqref{gap:eq:diagonal11} is equivalent to
\begin{equation}\label{gap:eq:harmonic-convolution}
 \sum_{h\geq1}\frac{z^h H_h(u)}{h+2u-1}
       =\frac z2\Phi(z,1,u)^2.
\end{equation}
This last elementary convolution identity can be proved directly by
partial fractions in the coefficient of $z^{h-1}$ in the square.
That observation clarifies the strength of the claim: the result is
an exact complete reduction, but the first-order instance already has
an elementary coefficient explanation.

\subsection{The half-shift diagonal in classical polylogarithms}

Put
\[
 \chi_s(x)=\frac{\operatorname{Li}_s(x)-\operatorname{Li}_s(-x)}2.
\]
For $z=x^2$ with $|x|<1$,
\[
 \Phi(z,s,\tfrac12)=\frac{2^s}{x}\chi_s(x),\qquad
 C_{2r}(\tfrac12)=2(2^{2r}-1)\zeta(2r),\qquad
 C_{2r+1}(\tfrac12)=0.
\]
The dependence on the chosen square root cancels in all products.
Equation \eqref{gap:eq:diagonal} gives the all-order identity
\begin{align}
 \mathcal M_{a,a}(z;\tfrac12,\tfrac12)
 =(-1)^a\Bigg[
 &2\sum_{r=1}^{\lfloor a/2\rfloor}
   (2^{2r}-1)\binom{2a-2r-1}{a-1}\zeta(2r)
                         \operatorname{Li}_{2a-2r}(z)\notag\\
 &-2^{2a-1}\chi_a(\sqrt z)^2\Bigg].
 \label{gap:eq:half-all}
\end{align}
The first four cases are
\begin{align}
 \mathcal M_{1,1}(z;\tfrac12,\tfrac12)
   &=2\operatorname{artanh}^2\sqrt z,\label{gap:eq:half11}\\
 \mathcal M_{2,2}(z;\tfrac12,\tfrac12)
   &=\pi^2\operatorname{Li}_2(z)-8\chi_2(\sqrt z)^2,
      \label{gap:eq:half22}\\
 \mathcal M_{3,3}(z;\tfrac12,\tfrac12)
   &=32\chi_3(\sqrt z)^2-3\pi^2\operatorname{Li}_4(z),
      \label{gap:eq:half33}\\
 \mathcal M_{4,4}(z;\tfrac12,\tfrac12)
   &=10\pi^2\operatorname{Li}_6(z)
       +\frac{\pi^4}{3}\operatorname{Li}_4(z)
       -128\chi_4(\sqrt z)^2.
      \label{gap:eq:half44}
\end{align}
Here each left side is a two-index series at a variable
argument; the displayed equality is a reduction of that function.
At individual special points there may be additional elementary
stuffle explanations, which should be recorded rather than presented
as new arithmetic phenomena.

\subsection{A classical antiderivative and an exact alternating moment}

At the half shift,
\[
 H_h(\tfrac12)=2H_{2h}-H_h.
\]
Hence \eqref{gap:eq:half11} says
\begin{equation}\label{gap:eq:harmonic-half}
 \sum_{h\geq1}\frac{(2H_{2h}-H_h)z^h}{h}
     =2\operatorname{artanh}^2\sqrt z.
\end{equation}
Integrating once against $dz/z$ is a useful explicit antiderivative.
For $0<z<1$ put \mbox{$x=(1-\sqrt z)/(1+\sqrt z)$}. Then
\begin{align}
 \sum_{h\geq1}\frac{(2H_{2h}-H_h)z^h}{h^2}
 ={}&\frac72\zeta(3)
   -\log^2x\log\frac{1+x}{1-x}\notag\\
   &+4\log x\,\chi_2(x)-4\chi_3(x).
 \label{gap:eq:primitive}
\end{align}
To prove it, write the left side as
$4\int_0^{\sqrt z}\operatorname{artanh}^2t\,dt/t$ and substitute
$x=(1-t)/(1+t)$. An antiderivative of
$\log^2x((1-x)^{-1}+(1+x)^{-1})$ is
\[
 \log^2x\log\frac{1+x}{1-x}
 -4\log x\,\chi_2(x)+4\chi_3(x).
\]
At $x=1$ its value is $4\chi_3(1)=7\zeta(3)/2$, fixing the constant.
The real interval formulation avoids any hidden logarithm branches;
continuation elsewhere is along the corresponding common path.

At the alternating endpoint one obtains
\begin{equation}\label{gap:eq:alternating-primitive}
 \sum_{h\geq1}\frac{(-1)^h(2H_{2h}-H_h)}{h^2}
       =\frac72\zeta(3)-2\pi G,
 \qquad G=\beta(2).
\end{equation}
An entirely real proof follows from
$-4\int_0^1\arctan^2t\,dt/t$; alternatively continue
\eqref{gap:eq:primitive} to $z=-1$, where $x=-i$ and the imaginary
$\pi^3$ terms cancel. The series is absolutely convergent.
This is a checkable example in the existing Euler-sum vocabulary,
not a claim that this classical type of special value is new.

\subsection{Rational shifts and finite root-of-unity filters}

For integers $1\leq p\leq q$, put $\omega=e^{2\pi i/q}$ and choose
$x$ with $x^q=z$. The residue-class filter gives
\begin{equation}\label{gap:eq:rational-filter}
 \Phi(z,s,p/q)=q^{s-1}x^{-p}
    \sum_{j=0}^{q-1}\omega^{-pj}
                    \operatorname{Li}_s(\omega^j x).
\end{equation}
For integer $s\geq1$ this follows directly by expanding the
polylogarithms and projecting onto exponents congruent to $p$ modulo
$q$. For larger positive rational shifts, first use
\[
 \Phi(z,s,q_0+N)=z^{-N}
 \left(\Phi(z,s,q_0)-\sum_{j=0}^{N-1}\frac{z^j}{(j+q_0)^s}\right).
\]
All apparent singularities at $z=0$ are removable. Thus rational
instances of \eqref{gap:eq:all} and \eqref{gap:eq:diagonal} reduce to
finite expressions involving ordinary polylogarithms at algebraically
related arguments and elementary cotangent derivatives. No extra
depth is introduced by evaluating the single Lerch functions.

\subsection{An independent integral check and reproducibility}

For $a,b\geq1$, $\operatorname{Re}u,\operatorname{Re}v>0$ and
$|z|<1$,
\begin{equation}\label{gap:eq:integral}
 \mathcal M_{a,b}(z;u,v)
 =\frac z{(a-1)!}\int_0^1
       \frac{t^u(-\log t)^{a-1}\Phi(t,b,v)}{1-zt}\,dt.
\end{equation}
Insert the Mellin integral for $(n+u)^{-a}$ and sum the geometric
series in $h=n-m$. This proves the formula by absolute convergence.
At a fixed $|z|=1$, $z\ne1$, its denominator stays away from zero on
$[0,1]$; the possible logarithmic singularity of $\Phi(t,1,v)$ at
$t=1$ is integrable. The boundary formula follows by dominated
convergence and the preceding summation-by-parts argument.

The supplied script \texttt{check\_shifted\_gap.py} evaluates this
integral after $t=e^{-y}$, using finite root-of-unity filters for
rational inner shifts. It compares that independent quadrature with
\eqref{gap:eq:all} for unequal shifts and complex colors, and checks
the first three half-diagonal formulas. At 60 decimal working digits,
the largest observed absolute residual among the four general
reflection checks is $8.33\times10^{-60}$. These checks are numerical
diagnostics, not rigorous error enclosures; the proofs are the
bilateral splitting and partial-fraction arguments above.

\subsection{Precise continuation questions}

\begin{enumerate}
\item Compare \eqref{gap:eq:all} term by term with the depth-two
specialization of Xu's shifted cyclotomic symmetry formulas, recording
the increasing/decreasing index conversion and color normalization.
The existence of a broad earlier theorem should remain explicit.
\item Determine which non-diagonal parameter specializations make the
reflected double equal to, or otherwise independently reducible from,
the original. Complementary shifts with equal exponents give the
complete reduction proved here.
\item Develop closed primitives of the complementary diagonal for
$a\geq2$ in a specified classical or multiple-polylogarithm basis.
The half-shift first-order primitive \eqref{gap:eq:primitive} supplies
a nontrivial normalization check.
\item Study noninteger spectral indices without differentiating the
integer-index partial-fraction formula as though its summation range
were analytic. Such an extension needs a separate contour or Mellin
argument; the integer formula alone does not justify it.
\item Continue in the shifts beyond the safe strip while tracking the
finitely many terms crossed when a shifted lattice point passes zero.
This should be a meromorphic identity with explicit finite
corrections, not an unqualified replacement of ordinary sums by
principal values.
\end{enumerate}
